\chapter{Breast Uplift}

\section*{Definition and Overview}
A breast uplift, medically termed mastopexy, is a cosmetic surgical procedure that raises and reshapes sagging breasts by removing excess skin and repositioning the breast tissue and the nipple--areola complex to a higher, firmer position. Unlike breast reduction, a breast uplift generally aims to preserve or maintain breast volume rather than reduce it. The procedure is usually performed under general anaesthesia, and the pattern of incisions depends on the degree of ptosis (sagging); some women require a scar limited to the areola, while others need a vertical scar or an anchor-shaped scar extending into the breast crease. A breast uplift may be combined with breast augmentation using implants when there is coexisting loss of volume, and it can also be performed after weight loss surgery to remove redundant skin. The primary goal is aesthetic: restoring breast shape, position, and projection while minimising scarring and preserving nipple sensation and the capacity to breastfeed where possible.

\section*{Epidemiology}
Mastopexy is one of the more commonly requested cosmetic breast procedures worldwide, and its popularity has grown steadily with the rise of body-contouring surgery after major weight loss. Women most often seek uplift surgery between their 30s and 50s, frequently after pregnancy, breastfeeding, or significant weight loss. Women with naturally large or soft breasts, thin or stretchy skin, and those who have experienced repeated weight fluctuations are particularly prone to drooping. The procedure is performed almost exclusively in women, although gynaecomastia surgery in men addresses a different problem and is not a mastopexy. Publicly funded surgery for breast uplift is limited in most healthcare systems, because the operation is regarded as cosmetic in the majority of cases; exceptions are occasionally made when marked ptosis causes physical symptoms such as rash, skin breakdown, or posture-related pain.

\section*{Aetiology and Pathophysiology}
Breast ptosis results from the combined effects of skin laxity, loss of skin elasticity, stretching of the supporting (Cooper's) ligaments, and reduction in the volume or fullness of glandular and fatty tissue.

\begin{itemize}
    \item \textbf{Ageing:} collagen and elastin fibres in the skin weaken and thin with age, so the breast skin loses its ability to resist the downward pull of gravity
    \item \textbf{Pregnancy and breastfeeding:} the breast enlarges and then involutes, stretching the skin and leaving glandular tissue less full, which produces a characteristic flattened, "deflated" appearance
    \item \textbf{Weight loss:} loss of subcutaneous fat reduces breast volume, leaving looser skin that hangs without firm support; this is especially marked after bariatric surgery
    \item \textbf{Gravity:} the constant downward force of gravity acts throughout life and is more pronounced in heavier, larger breasts
    \item \textbf{Genetic factors and body habitus:} some women inherit naturally soft, less dense breast tissue or thin, stretchy skin
    \item \textbf{Smoking:} smoking damages dermal elastic fibres and impairs blood supply to the skin, hastening sagging and increasing surgical risks
    \item \textbf{Hormonal changes:} fluctuations during the menstrual cycle, menopause, and oral contraceptive use can alter breast volume and fullness over time
\end{itemize}

The degree of ptosis is commonly graded using the Regnault classification, which relates the position of the nipple to the inframammary fold (the crease beneath the breast); the grading guides the choice of incision pattern.

\section*{Clinical Features}
Sagging of the breasts is a normal anatomical variation rather than a disease, but women present for assessment with a combination of recognisable features.

\begin{itemize}
    \item The nipple points downward or lies at or below the level of the inframammary fold
    \item The upper part of the breast appears empty, flat, or deficient in fullness
    \item Loose, redundant skin around the lower breast and areola
    \item The areola appears stretched, widened, or enlarged
    \item The breast may feel "deflated" after pregnancy or weight loss
    \item Difficulty fitting bras and clothing, chafing, and moisture or rash beneath the breast
    \item Self-consciousness about appearance, which may affect body image, confidence, and quality of life
\end{itemize}

\section*{Differential Diagnosis}
\begin{itemize}
    \item \textbf{Normal age-related ptosis:} the commonest reason women request uplift, and not a medical problem
    \item \textbf{Pseudoptosis:} the lower breast droops while the nipple remains above the fold, usually from loss of volume rather than skin laxity
    \item \textbf{Breast cancer or a breast lump:} a new hard lump, tethering, dimpling of the skin, or a change in breast shape or nipple position warrants urgent investigation
    \item \textbf{Mastitis or breast abscess:} a swollen, painful, reddened breast, particularly during breastfeeding, with systemic features in the case of abscess
    \item \textbf{Inflammatory breast changes:} diffuse redness, swelling, and firmness of the breast, which may indicate an aggressive inflammatory cancer
    \item \textbf{Benign breast conditions:} cysts, fibroadenomas, or fat necrosis that alter the contour or size of one breast
\end{itemize}

\section*{When to Seek Medical Care}
A breast uplift is not a medical treatment, so the main reasons to consult a doctor are to confirm that sagging is a normal variation, to investigate any concerning features, and to obtain a balanced assessment of whether surgery is appropriate. Any new, unexplained breast change should always be discussed with a doctor, regardless of cosmetic concerns.

\textbf{Call 000 (or local emergency services) for:}
\begin{itemize}
    \item A new breast lump that is hard, fixed, or rapidly growing
    \item Skin dimpling, puckering, or orange-peel changes over the breast
    \item Bloody or spontaneous discharge from the nipple
    \item Sudden redness, swelling, heat, or fever suggesting infection
    \item A new inversion or pulling inward of the nipple
\end{itemize}

\section*{Diagnosis and Investigations}
\begin{itemize}
    \item \textbf{History:} questions about weight changes, pregnancies, breastfeeding history, smoking, medications, and family history of breast cancer
    \item \textbf{Examination:} clinical assessment of the degree of ptosis (Regnault grading), skin quality, breast volume, symmetry, and the position of the nipple relative to the inframammary fold
    \item \textbf{Photographic documentation:} standardised photographs taken before surgery to plan the procedure and allow comparison of results
    \item \textbf{Breast imaging:} mammography and/or ultrasound in women aged over 40, in those with a family history of breast cancer, or where any lump or concerning feature is found, to exclude malignancy
    \item \textbf{General health assessment:} a preoperative review of fitness for anaesthesia, including blood pressure, blood counts, and cardiovascular risk where indicated
    \item \textbf{Counselling and informed consent:} a detailed discussion of realistic expectations, incision patterns, recovery, costs, and potential complications
\end{itemize}

\section*{Management}
\begin{itemize}
    \item \textbf{Non-surgical measures:} a well-fitted, supportive bra can improve comfort and perceived shape; weight stability, skin care, and exercise help general skin quality but cannot lift breast tissue
    \item \textbf{Periareolar mastopexy:} a crescent or circular incision around the areola, suitable for mild ptosis
    \item \textbf{Vertical (lollipop) mastopexy:} an incision around the areola with a vertical limb to the breast crease, used for moderate ptosis
    \item \textbf{Inverted-T or anchor mastopexy:} an incision around the areola, vertically down the breast, and along the fold, used for more significant ptosis and excess skin
    \item \textbf{Mastopexy with implants:} a combined procedure that restores both shape and volume in women who have lost fullness
    \item \textbf{Post-operative care:} dressings, a surgical or supportive bra, drainage when required, and avoidance of heavy lifting and strenuous exercise for several weeks
    \item \textbf{Surgeon counselling:} discussion of realistic expectations, scar patterns and their permanence, and the natural progression of ageing after surgery
\end{itemize}

\section*{Complications}
\begin{itemize}
    \item Bleeding, bruising, or a fluid collection (seroma or haematoma) under the skin, occasionally needing drainage
    \item Infection or poor wound healing, particularly in smokers
    \item Changes in nipple sensation, which may be temporary or permanent
    \item Difficulty breastfeeding in the future if the milk ducts or nipple areola complex are affected
    \item Scarring that is prominent, widened, or slow to fade
    \item Asymmetry, unsatisfactory shape, or "bottoming out" with recurrent drooping
    \item Rarely, loss of part of the nipple or areola if the blood supply to the skin is damaged
    \item If implants are used, the additional risks of implant rupture, infection, capsular contracture, and the need for further surgery
    \item General risks of anaesthesia, including deep vein thrombosis, particularly in longer procedures
\end{itemize}

\section*{Prognosis and Outcomes}
Most women are satisfied with the results of mastopexy and report a firmer, more youthful breast shape and improved confidence. Scars typically fade and soften over 6--12 months but remain visible. The result is not permanent: gravity, ageing, pregnancy, breastfeeding, and weight fluctuation continue to act on the breasts, and some women later consider revision surgery. Complications such as infection and poor healing are significantly more common in smokers, so stopping smoking well before surgery is strongly advised. Careful selection of a qualified, experienced surgeon, adherence to post-operative instructions, and realistic expectations are the most important factors in achieving a good outcome.

\section*{Prevention and Self-Care}
\begin{itemize}
    \item Avoid large, rapid weight changes, which stretch the skin and contribute to sagging
    \item Wear a supportive, well-fitted bra, particularly during pregnancy, breastfeeding, and exercise
    \item Stop smoking, since smoking reduces skin elasticity and impairs wound healing
    \item Protect the breast skin from sun damage and keep it well moisturised
    \item Maintain a healthy weight and overall fitness
    \item Understand that creams, massage, and exercise cannot lift breast tissue or reverse stretched skin
    \item If considering surgery, choose a qualified and experienced surgeon, ask about revision policy, and attend all follow-up appointments
\end{itemize}

\section*{Sources and Further Reading}
The following organisations provide reliable information on cosmetic breast surgery, mastopexy, and breast health for patients and clinicians.

\begin{itemize}
    \item NHS. \textit{Cosmetic breast surgery, including breast uplift, and what to consider before surgery.} https://www.nhs.uk/conditions/cosmetic-surgery/
    \item Mayo Clinic. \textit{Breast lift (mastopexy): what you can expect and recovery.} https://www.mayoclinic.org/
    \item MSD Manual. \textit{Breast disorders and cosmetic breast surgery overview.} https://www.msdmanuals.com/
    \item MedlinePlus. \textit{Breast lift surgery and cosmetic breast procedures.} https://medlineplus.gov/breastsurgery.html
    \item American Society of Plastic Surgeons. \textit{Breast lift surgery information and patient safety resources.} https://www.plasticsurgery.org/
    \item National Health and Medical Research Council. \textit{Consumer health and clinical guidance.} https://www.nhmrc.gov.au/
\end{itemize}
